\documentclass[11pt,a4paper]{article}
\usepackage[T1]{fontenc}
\usepackage[utf8]{inputenc}
\usepackage{amsmath,amssymb,xcolor}
\usepackage[margin=25mm]{geometry}
\usepackage[colorlinks=true,urlcolor=blue]{hyperref}
\definecolor{nexorablue}{HTML}{2459B5}
\setlength{\parindent}{0pt}
\setlength{\parskip}{7pt}
\setlength{\emergencystretch}{2em}
\renewcommand{\familydefault}{\sfdefault}
\begin{document}
\makeatletter
\def\@oddfoot{\footnotesize Download for free at \url{https://openstax.org/details/books/university-physics-volume-2}.\hfil\thepage}
\let\@evenfoot\@oddfoot
\makeatother
{\small\bfseries\color{nexorablue} NEXORA CPGE}\par
{\small MP / OST-SPE-13-01 / Énoncé}\par
{\LARGE\bfseries Bobine carrée dans un champ décroissant}\par
\vspace{0.25em}\hrule\vspace{0.6em}

Une bobine carrée de côté $0{,}25\,\mathrm m$ possède $200$ spires et présente une résistance de $5{,}0\,\mathrm\Omega$. Elle est fixe dans un champ uniforme normal à son plan, dont la norme diminue à $0{,}040\,\mathrm{T\,s^{-1}}$. Déterminer les valeurs absolues de la force électromotrice et du courant. On néglige l’auto-induction.

\bigskip
{\footnotesize Adaptation française par Nexora CPGE. Source : Samuel J. Ling, Jeff Sanny et William Moebs (OpenStax, Rice University),
\href{https://d3bxy9euw4e147.cloudfront.net/oscms-prod/media/documents/UniversityPhysicsVolume2-OP.pdf\#page=595}{\textit{University Physics Volume 2, édition 2018}}, § 13.1, exemple 13.1, p. 585.
\href{https://creativecommons.org/licenses/by/4.0/}{CC BY 4.0}.}
\end{document}
